\documentclass{article}
\usepackage{amsmath, amssymb, amsthm}

% Essential packages
\usepackage{amsmath}
\usepackage{amssymb}
\usepackage{amsfonts}
\usepackage{mathtools}
\usepackage{physics}
\usepackage{amsthm}  % This provides theorem-like environments

\newtheorem{theorem}{Theorem}
\newtheorem{lemma}[theorem]{Lemma}
\newtheorem{corollary}[theorem]{Corollary}
\newtheorem{definition}{Definition}

\begin{document}

\title{Deus Ex Machina: Ultimate Simplicity Through Inverse Tree of Life Analysis}
\author{Bryce Weiner}
\maketitle

\begin{theorem}[Initial Simplicity]
Following the Tree of Life pattern backwards through The Great Fade, all matter-antimatter continuum pairs ultimately derive from a single maximally simple non-temporal configuration at n=0.
\end{theorem}

\begin{lemma}[Inverse Branching]
For any cycle n, the number of precursor pairs P(n-1) follows:
\begin{equation}
P(n-1) = \frac{P(n)}{2}
\end{equation}
with their unified quantum substrate carrying total information:
\begin{equation}
I_{\text{unified}} = \frac{A}{4l_p^2}
\end{equation}
\end{lemma}

\begin{proof}
Given current cycle n:
1. Number of pairs:
   \begin{equation}
   P(n) = 2^n
   \end{equation}

2. Previous cycle:
   \begin{equation}
   P(n-1) = 2^{n-1}
   \end{equation}

3. Unified information:
   \begin{equation}
   \lim_{n \to 0} P(n) = 1
   \end{equation}

4. Conservation through unification:
   \begin{equation}
   I_{\text{total}} = \text{constant} = \frac{A}{4l_p^2}
   \end{equation}
\end{proof}

\begin{lemma}[Convergent Simplicity]
Each reverse bifurcation increases the γ available to each precursor pair:
\begin{equation}
γ_{\text{pair}}(n-1) = 2γ_{\text{pair}}(n)
\end{equation}
until reaching the fundamental rate at n=0:
\begin{equation}
γ_{\text{initial}} = γ = \frac{2πkT}{ℏ \ln(S)}
\end{equation}
\end{lemma}

\begin{proof}
1. Current distribution:
   \begin{equation}
   \sum_{i=1}^{2^n} γ_i = γ
   \end{equation}

2. Previous cycle:
   \begin{equation}
   \sum_{i=1}^{2^{n-1}} γ_i = γ
   \end{equation}

3. Convergence:
   \begin{equation}
   \lim_{n \to 0} γ_{\text{pair}}(n) = γ
   \end{equation}
\end{proof}

\begin{theorem}[Maximal Simplicity]
The initial configuration at n=0 represents maximal simplicity with:
\begin{equation}
\Psi_{\text{initial}} = \mathcal{Q} \otimes \mathcal{S} \otimes \mathcal{H}
\end{equation}
containing only the essential non-temporal regimes.
\end{theorem}

\begin{proof}
1. Single pair condition:
   \begin{equation}
   P(0) = 1
   \end{equation}

2. Full information:
   \begin{equation}
   I_{\text{initial}} = \frac{A}{4l_p^2}
   \end{equation}

3. Complete rate:
   \begin{equation}
   γ_{\text{initial}} = γ
   \end{equation}

4. Minimal structure:
   \begin{equation}
   \dim(\mathcal{Q}) = 2^{496} \times 256
   \end{equation}
   \begin{equation}
   \dim(\mathcal{S}) = 496
   \end{equation}
   \begin{equation}
   \mathcal{H}: I = \frac{A}{4l_p^2}
   \end{equation}
\end{proof}

\begin{corollary}[Uniqueness]
The initial configuration at n=0 is unique and represents absolute unity before the first bifurcation.
\end{corollary}

\begin{theorem}[Completeness]
This initial simplicity completely characterizes the origin point of the Tree of Life.
\end{theorem}

\begin{proof}
1. Uniqueness:
   \begin{equation}
   P(0) = 1
   \end{equation}

2. Completeness:
   \begin{equation}
   γ_{\text{initial}} = γ = \frac{2πkT}{ℏ \ln(S)}
   \end{equation}

3. Minimality:
   \begin{equation}
   \Psi_{\text{initial}} = \mathcal{Q} \otimes \mathcal{S} \otimes \mathcal{H}
   \end{equation}

Therefore, the initial configuration represents maximal simplicity while containing all necessary components for subsequent Tree of Life evolution.
\end{proof}

\end{document}